\section{Introduction}\label{sec:intro}

A microservice application is a graph of independently deployed services that call each other over the network and share
hosts, caches and databases. When one component fails, the failure crosses component boundaries: an exhausted CPU on a node slows
every service it hosts, a failing payment service returns errors to the checkout service that calls it, and the operator sees
alerts far from the cause \citep{avizienis2004taxonomy,soldani2022anomaly}. Root cause analysis (RCA) must therefore answer three
questions from the incident's telemetry---metrics, logs, traces, Kubernetes events and the service topology: \emph{which} entity is
the root cause, \emph{what} fault mechanism acts there, and \emph{how} the fault propagated to the symptoms
\citep{zhang2025survey,fu2025survey}. Recent benchmarks score all three. RCA100 \citep{cai2026benchmark} grades entity localisation,
fault identification over 28 fault types and the reasoning process along a typed causal chain with numeric evidence checkpoints.

Two families of methods address the task. Statistical and causal methods rank candidate components from metric changes, trace
structure or learned causal graphs \citep{wu2020microrca,ikram2022rcd,pham2024baro,li2022circa,yu2023nezha,zheng2024mulan}. They are
cheap and reproducible, but they usually return a ranked list of components without a fault type or an explanation. Agents built on
large language models (LLMs) read the telemetry through tools and write a diagnosis in natural language
\citep{chen2024rcacopilot,wang2024rcagent,xu2025openrca,pei2025flowofaction,zhang2024mabc}. They produce types and explanations,
but a trajectory-level study of such agents finds systematic reasoning failures---untested baselines, ignored timing, reversed
caller / callee semantics, anchoring on the first hypothesis \citep{lu2026beyond,riddell2026stalled}---and the strongest agents
consume on the order of $10^6$ tokens per incident \citep{lu2026beyond,luo2026pos}. Agents that maintain an explicit belief state
\citep{luo2026pos,luo2026graphofstates} still ask the LLM to \emph{write} that belief.

This paper takes the opposite position: the belief over root causes can be \emph{computed} from calibrated evidence, and the
language model, if used at all, should be bounded to the parts that need language. We present \gama{} (Graph-aware, Multimodal,
Abductive RCA). Every continuous signal of every entity is scored against the incident's own history and its peers, and the score
is turned into a p-value by an empirical null cut from that history (pseudo-windows); events, status codes, failing spans and log
text become categorical observations. A Bayesian model over (entity, fault type) hypotheses scores these observations with
fingerprints that state, for each fault type, which signals are key, supporting, possible consequences or contradictions. Symptoms
that a hypothesis can explain along the topology are explained away; the propagation strength is a per-incident nuisance parameter
estimated by expectation--maximisation; and a root may not start after the symptoms it explains. The decision is two-stage (entity,
then type), and the causal chain is computed along the topology, citing the window values of every step. An optional LLM may read
log text chosen by expected information gain, choose a type from the statistical shortlist with a verbatim quote, and narrate the
chain.

We evaluate \gama{} on the 103 cases of RCA100 under a pre-registered protocol (hypotheses, statistics and every amendment
recorded before the corresponding run) and release every prompt and output. The results support three statements. First, the
statistical core alone reaches a Final score of 62.9, on a CPU and without labels, 20.2 points above an LLM that decides from the
same calibrated evidence and 44.1 points above an LLM that reads the alert only. Second, the bounded LLM lowers the decision score
by 3.7 (Qwen2.5-7B) and 3.5 points (gpt-oss-120b): when the telemetry cannot separate two fault types, the LLM tends to pick the
more specific one and is more often wrong than right. Third, scale improves the \emph{explanation}: the narrator of the larger model
hits 2.7 times as many step-anchored evidence checkpoints as the 7B narrator, close to the computed chain. We then transfer the
method zero-shot to AIOps2025 \citep{cai2026benchmark}, a different system with different fault types, changing only
dataset knowledge (an adapter and a fingerprint table written from the published fault definitions).

\paragraph{Contributions}
\begin{enumerate}
\item \textbf{A computed abductive belief for RCA.} A label-free likelihood model over (entity, fault type) hypotheses with
fingerprint roles, topology explain-away, per-incident propagation strength (EM) and temporal precedence (\cref{sec:method}).
\item \textbf{Calibrated multimodal evidence.} Shift scores against the incident's own history and peers, calibrated on
pseudo-windows, together with trace-origin attribution, direction-aware load explain-away, mechanism-naming text and cross-modal
corroboration as separable modules (\Mod{1}--\Mod{6}).
\item \textbf{A pre-registered study of what an LLM adds.} Seven arms on the same 103 cases, from an LLM that sees only the alert
to the computed belief alone, two backbones (7B local, 120B through an API) and hypotheses fixed in advance
(\cref{sec:setup,sec:results}): the LLM does not improve the decision at either scale, while scale improves the explanation.
\item \textbf{A zero-shot transfer protocol} to an unseen system, pre-registered with an implementation audit, and an open
implementation with equivalence and synthetic tests (\cref{sec:aiops}).
\end{enumerate}

The paper is organised as follows. \Cref{sec:related} reviews related work, \cref{sec:problem} states the problem and the metrics,
\cref{sec:method} describes the method, \cref{sec:setup} the experimental protocol, \cref{sec:results} the results,
\cref{sec:discussion} discusses implications and threats to validity, and \cref{sec:conclusion} concludes.
